% ===========================================================================
% IV. MARL Framework
% ===========================================================================
\section{Multi-Agent Reinforcement Learning Framework}
\label{sec:marl}

\subsection{Environment Design}
\label{sec:env}

The environment is implemented as a PettingZoo AEC environment~\cite{terry2021pettingzoo}
with $N = 20$ agents.  At each step, the environment selects one agent in
round-robin order; that agent observes the current state, selects an action,
and receives an immediate scalar reward.  An \emph{episode} consists of
$T = 100$ steps shared across all agents (i.e., each agent acts approximately
5 times per episode).

\subsubsection{Observation Space}
Each agent $i$ observes a 14-dimensional vector:
\begin{equation}
  o_i = [\mu_i,\ w_i,\ s_i,\ \alpha_i,\ \beta_i,\ \mu_j,\ w_j,\
         r_i,\ d_i,\ d^{\text{lost}}_i,\ t,\ n^{\text{sybil}}_i,\
         B_{\text{norm}},\ g_i],
  \label{eq:obs}
\end{equation}
where $\mu_i$ is agent $i$'s own score, $w_i$ its CI width, $s_i$ its
stake, $\alpha_i$ and $\beta_i$ its Beta parameters, $\mu_j$ and $w_j$
are the current target's score and CI width, $r_i$ is the number of ratings
submitted this episode, $d_i$ and $d^{\text{lost}}_i$ are disputes received
and lost, $t \in [0,1]$ is normalized episode time, $n^{\text{sybil}}_i$
is the agent's Sybil count, $B_{\text{norm}}$ is the normalized attack-block
count, and $g_i \in \{0,1\}$ is gate eligibility.  All continuous fields
are clipped to $[-5, 5]$.

\subsubsection{Action Space}
The 12-dimensional discrete action space is:
\begin{enumerate}[label=\textbf{A\arabic*.}, leftmargin=*, noitemsep, start=0]
  \item No-operation (abstain)
  \item Submit positive honest rating
  \item Submit negative honest rating
  \item Submit positive dishonest rating
  \item Submit negative dishonest rating
  \item Create Sybil identity
  \item File dispute
  \item Attempt self-rating (attack)
  \item Attempt admin escalation (attack)
  \item Tamper with evidence (attack)
  \item Attempt gate bypass (attack)
  \item Attempt provenance replay (attack)
\end{enumerate}
Actions 7--11 are the five explicit attack-class actions corresponding to
the five programmable attack vectors in the blockchain companion
system~\cite{almaayn2024unified}; actions 3--5 implement the remaining
attack vectors (dishonest rating, Sybil) through non-attack actions with
adversarial intent.

\subsubsection{Reward Function}
The reward for agent $i$ at step $t$ is:
\begin{equation}
  r_i^t = R_{\text{base}}(a_i^t) + R_{\text{defense}}(a_i^t) + R_{\text{dividend}}^t,
  \label{eq:reward}
\end{equation}
where:
\begin{align}
  R_{\text{base}}(a) &= \begin{cases}
    +0.1 & a \in \{1,2\}\ (\text{honest rating}) \\
    -0.5 & a \in \{3,4\}\ (\text{dishonest, detected}) \\
    +b   & a \in \{3,4\}\ (\text{dishonest, undetected}) \\
    0.0  & a = 0
  \end{cases} \label{eq:base_reward} \\
  R_{\text{defense}}(a) &= \begin{cases}
    -5.0 & a \in \{7, 8, 11\}\ (\text{deterministic block}) \\
    -2.0 & a = 9 \wedge \text{detected}\ (p=0.80) \\
    -3.0 & a = 10\ (\text{gate bypass}) \\
    0.0  & \text{otherwise}
  \end{cases} \label{eq:defense_reward} \\
  R_{\text{dividend}}^t &= +0.1 \cdot \mathbb{1}[\text{attack blocked this step}].
  \label{eq:dividend}
\end{align}
The defense dividend $R_{\text{dividend}}^t$ is received by \emph{all} honest
agents when any attack is blocked in the current step, creating a collective
incentive to indirectly reinforce the defense layer.  Adversarial agents
have an additional reward bonus $b \in [0.5, 2.0]$ for successful dishonest
actions, as specified per configuration in Table~\ref{tab:configs}.

\subsection{MAPPO Training}
\label{sec:mappo}

We train agents using MAPPO with a shared actor-critic
network~\cite{yu2021surprising}.  The network architecture is:
\begin{equation}
  f_\theta(o) = \text{ReLU}(W_2\,\text{ReLU}(W_1 o + b_1) + b_2),
  \label{eq:network}
\end{equation}
where $W_1 \in \mathbb{R}^{256 \times 14}$ and $W_2 \in \mathbb{R}^{128
\times 256}$, followed by separate linear actor and critic heads.  Sharing
parameters across all agents means that the policy represents a general
strategy conditioned only on the agent's own observation; this is appropriate
because all agents are structurally symmetric (they differ only in their
current state, not their capabilities).

\subsubsection{PPO Update}
At the end of each episode, the rollout buffer contains all transitions from
all agents.  GAE advantages are computed per-agent:
\begin{equation}
  \hat{A}_t = \sum_{l=0}^{\infty}(\gamma\lambda)^l \delta_{t+l},\quad
  \delta_t = r_t + \gamma V(o_{t+1}) - V(o_t),
  \label{eq:gae}
\end{equation}
with $\gamma = 0.99$ and $\lambda = 0.95$.  All transitions are pooled into
a single minibatch for the PPO clipped surrogate update:
\begin{equation}
  \mathcal{L}_{\text{clip}} = \mathbb{E}_t\Bigl[
    \min\bigl(r_t(\theta)\hat{A}_t,\
    \text{clip}(r_t(\theta), 1\!-\!\epsilon, 1\!+\!\epsilon)\hat{A}_t\bigr)
  \Bigr],
  \label{eq:ppo_clip}
\end{equation}
where $r_t(\theta) = \pi_\theta(a_t|o_t)/\pi_{\theta_{\text{old}}}(a_t|o_t)$
and $\epsilon = 0.2$.  The full loss is:
\begin{equation}
  \mathcal{L} = -\mathcal{L}_{\text{clip}} + c_v \mathcal{L}_{\text{VF}} - c_e H[\pi_\theta],
  \label{eq:full_loss}
\end{equation}
with value coefficient $c_v = 0.5$, entropy coefficient $c_e = 0.05$ (higher
than the typical 0.01 to encourage early exploration of all 12 actions), and
value function loss $\mathcal{L}_{\text{VF}} = \text{MSE}(V(o_t), R_t)$.

\subsubsection{Learning Rate Schedule}
The learning rate decays linearly from $3\times10^{-4}$ to $3\times10^{-5}$
over the training run:
\begin{equation}
  \eta_t = 3\times10^{-4} \cdot \max\!\left(0.1,\ 1 - \frac{t}{T_{\max}}\right),
  \label{eq:lr_schedule}
\end{equation}
where $T_{\max}$ is the maximum episode count (\icode{max\_episodes\_extended}).
This schedule prevents over-fitting to early-episode dynamics and ensures
convergence at long horizons.

\subsection{Experimental Configurations}
\label{sec:configs}

Table~\ref{tab:configs} summarizes the 11 experimental configurations.
Each configuration specifies the adversarial population size, the subset of
enabled attack-class actions, and the adversarial reward bonus.  Configurations
C1--C5 establish baseline and mixed-population scenarios; C6--C10 isolate
individual attack vectors; C11 enables all 9 attack types simultaneously
with a colluding subgroup and Sybil capability.

\begin{table}[!t]
\centering
\caption{Experimental configurations.  $N_{\text{adv}}$ = adversarial agents
(out of 20). $b$ = adversarial reward bonus. ``Enabled'' = attack-class
actions enabled (all configs also permit actions 0--6).}
\label{tab:configs}
\footnotesize
\renewcommand{\arraystretch}{1.2}
\begin{tabular}{@{}C{0.5cm}L{2.4cm}C{0.8cm}C{0.6cm}L{2.5cm}@{}}
\toprule
\textbf{C\#} & \textbf{Name} & $N_{\text{adv}}$ & $b$ & \textbf{Enabled Attacks} \\
\midrule
1  & Baseline (all honest)   & 0  & 0.0 & None \\
2  & Mixed population        & 5  & 0.5 & None (A3--A6) \\
3  & Sybil                   & 5  & 0.5 & A5 \\
4  & Collusion ring          & 5  & 0.5 & A3,A4 (group) \\
5  & Adaptive (2$\times$)    & 10 & 2.0 & A3--A6 \\
6  & Self-rating             & 5  & 1.0 & A7 \\
7  & Admin escalation        & 3  & 1.0 & A8 \\
8  & Evidence tampering      & 5  & 1.0 & A9 \\
9  & Gate bypass             & 4  & 1.0 & A10 \\
10 & Provenance replay       & 3  & 1.0 & A11 \\
11 & Comprehensive           & 10 & 1.0 & A7--A11 + A3--A6 \\
\bottomrule
\end{tabular}
\end{table}

\subsection{Training Protocol}

Each configuration is trained with $S = 5$ independent random seeds (seeds
0--4).  Each seed runs for up to $T_{\max} = 10{,}000$ episodes, with early
stopping if the reward variance over the last 100 episodes falls below
$\sigma^2 = 0.05$.  Evaluation metrics are computed as the mean over the
final 20\% of training episodes (episodes 8{,}000--10{,}000), avoiding
contamination from early random-policy behavior.  All training runs use
\texttt{torch.set\_num\_threads(1)} with \icode{OMP\_NUM\_THREADS=1} and
\icode{MKL\_NUM\_THREADS=1} to prevent CPU thread contention during parallel
runs.  Wall-clock budget is 8 hours per seed.

The primary metrics are:
\begin{itemize}[noitemsep]
  \item \textbf{Honest \%}: Fraction of rating actions (A1--A4) in which the
        agent chose an honest action (A1 or A2).
  \item \textbf{Defense block rate}: $\text{blocked} / \text{attempted}$
        for all attack-class actions (A5--A11) in the tail.
  \item \textbf{Reputation accuracy}: Mean absolute error between agent score
        $\mu_i$ and true latent quality $q_i \sim \mathcal{U}[0,1]$.
\end{itemize}
